\clearpage
\renewcommand{\sectioncolor}{ember}
\section{AlphaFold\,3: exactly what it cracked, and what it left untouched}
\markboth{AlphaFold}{}

I do not want to undersell this. AF2 and AF3 are a genuine landmark. But the precise statement of
what was solved is what tells you where to work next.

\subsection{What AF3 is, mathematically}

\begin{evidence}[title={Abramson et al., \textit{Nature} \textbf{630}, 493 (May 2024)}]
AF2's structure module used \textbf{Invariant Point Attention} --- SE(3)-equivariant attention over
per-residue frames (a rotation and a translation per residue), so the prediction is independent of
how you orient the whole molecule.

\textbf{AF3 replaced that module with a diffusion model operating directly on atomic coordinates},
and extended coverage to nucleic acids, ligands, ions and modified residues. Equivariance is traded
for a generative denoiser plus data augmentation.
\end{evidence}

\noindent Note the irony, and it is a useful one: \textbf{AF3's structure module is already a
generative model of a distribution} --- it just gets collapsed to a single ranked output. The
machinery for producing a measure is partly there; what is missing is that the model is not trained
to make that measure the \emph{Boltzmann} measure.

\subsection{The honest ledger}

\begin{center}\small
\begin{tabular}{@{}>{\raggedright\arraybackslash}p{4.2cm}>{\raggedright\arraybackslash}p{5.9cm}>{\raggedright\arraybackslash}p{6.2cm}@{}}
\toprule
\rowcolor{ember!13}
\textbf{\color{ember}Left untouched} & \textbf{\color{ember}Why it matters} & \textbf{\color{ember}State of play}\\
\midrule
\textbf{Conformational ensembles} &
 Function lives in the ensemble: allostery, cryptic pockets, induced fit &
 MSA-subsampling hacks (Del Alamo et al., \textit{eLife} \textbf{11}, e75751, 2022); AlphaFlow and
 ensemble emulators are early and unvalidated\\
\textbf{Free energies and affinities} &
 AF3 returns a \emph{pose}, not $\Delta G$ &
 Affinity prediction remains the open problem; see §4\\
\textbf{Effect of a point mutation} &
 A single substitution barely changes the MSA, so the prediction barely moves &
 AF2 performs poorly on $\Delta\Delta G$ of stability --- repeatedly documented\\
\textbf{Intrinsic disorder} &
 $\sim$30\,\% of the human proteome &
 Low-pLDDT spaghetti: honest, but not usable\\
\rowcolor{ember!7}
\textbf{Kinetics} &
 Folding paths, transition states, residence time --- the actual pharmacological variable &
 Entirely outside the model\\
\textbf{Novel folds, shallow MSAs} &
 Performance degrades where there is no family to lean on &
 \textbf{The same failure mode flagged for AMR in Part VII:} a learned prior regresses the novel
 toward the common\\
\bottomrule
\end{tabular}
\end{center}

\begin{tier3}[title={The one-sentence statement}]
\begin{center}\itshape\large
AlphaFold solved a \emph{structure} problem.\\
Biology needs an \emph{ensemble-and-rate} problem solved.
\end{center}
\vspace{2pt}
These are different mathematical objects. Formally, what is missing is the Boltzmann measure $\mu$
and the \textbf{transfer operator} $\mathcal{P}_t$ that propagates it. A structure predictor returns
a point. Biology needs $(\mu,\mathcal{P}_t)$ --- and the second one has barely been started.
\end{tier3}

\begin{plain}
AlphaFold gave us a very good photograph of a machine. What we need is a film of it working, and a
measurement of how fast each part moves. A photograph is not a bad film --- it is a different kind of
object, and no amount of sharpening turns one into the other.
\end{plain}
